% !TEX program = xelatex
\documentclass[11pt,a4paper]{article}
\usepackage{fontspec}
\usepackage[margin=2.5cm]{geometry}
\usepackage[hidelinks,colorlinks=false]{hyperref}
\usepackage{parskip}
\pagestyle{empty}

\begin{document}

18 August 2026

Editorial Team\\
\textit{Scientific Reports}

\textbf{Re: Article submission, ``Evidence-Qualified Alignment of Pathology Foundation
Models with Morphologic, Molecular, and Outcome Axes in Prostate Cancer''}

Dear Editors,

Please consider our manuscript, ``Evidence-Qualified Alignment of Pathology Foundation
Models with Morphologic, Molecular, and Outcome Axes in Prostate Cancer,'' for publication
as an Article in \textit{Scientific Reports}.

Pathology foundation models can encode clinically meaningful information, but recovery of a
target from a representation does not establish that a downstream model uses that target or
that the relationship will transport across cohorts. We address this problem by comparing six
non-interchangeable prostate-cancer axes---grade/ISUP, tumor phenotype/content, PTEN loss,
androgen-receptor activity, SPOP mutation, and recurrence---across frozen CONCH and Virchow
representations. The study evaluates recoverability together with external transport,
conditional information, site uncertainty, encoder and sampling sensitivity, endpoint
fidelity, and a narrowly scoped internal functional-sensitivity analysis.

The results establish a deliberately qualified evidence hierarchy. Grade/ISUP and tumor
phenotype/content show the strongest representation evidence; PTEN and androgen-receptor
activity remain context-sensitive; SPOP is evaluated but unsupported in the frozen design;
and recurrence depends on endpoint and reference specification. Targeted removal of an
ISUP-correlated direction changes two internally locked biochemical-recurrence heads beyond
matched-random controls, but the study does not claim indispensable mechanism, external
functional transport, clinical increment, or improved clinical outcomes. We believe this
combination of multi-axis evidence, explicit claim boundaries, and auditable provenance is
well aligned with the broad methodological and biomedical readership of
\textit{Scientific Reports}.

The manuscript reports secondary analyses of deidentified data from previously published
public or access-governed repositories. No new participants were recruited, and no identifiable
participant information is reported. Original ethical approvals and consent or waivers remain
those reported by the source studies and repositories. The authors declare no competing
interests. The work was supported by NRF-2023R1A2C1006639 and in part by Gyeongsang National
University.

Aggregate source tables, provenance registries, analysis-code snapshots, and deterministic
figure-reproduction tests are available in the
\href{https://github.com/AiXLab-GNU/evidence-qualified-alignment-prostate-cancer}{public release repository} under the
immutable release tag \texttt{v1.0.4-submission}. Reviewer suggestions or exclusions and the
status of any prior discussion with a Scientific Reports Editorial Board Member will be
declared by the corresponding author in the online submission system.

Thank you for considering our manuscript.

Sincerely,

\textbf{Jin Hyun Kim}\\
Associate Professor\\
Gyeongsang National University\\
Jinju, Republic of Korea\\
Email: \href{mailto:jin.kim@gnu.ac.kr}{jin.kim@gnu.ac.kr}\\
ORCID: \href{https://orcid.org/0000-0002-2308-1638}{0000-0002-2308-1638}\\
Corresponding author

\end{document}
